% ============================================================
% CHAPTER 14 — Testing and Security
% ============================================================
\chapter{Testing and Security}\index{testing}\index{sicurezza}

\section*{What You Will Build}
A testing and security verification system for the entire stack:
\begin{itemize}
  \item AI-generated unit tests for backend and frontend
  \item Integration tests for API endpoints
  \item Automated OWASP Top~10 security analysis
  \item Confidence Tagging to decide what to review manually
  \item A reproducible quality pipeline
\end{itemize}

\textbf{Estimated time}: 60--90 minutes\\
\textbf{Prerequisite}: Complete full-stack application (Chapters~6--13)

\begin{dangerbox}[title={Security Is Not Optional}]
AI-generated code can contain security vulnerabilities --- not out of malice, but because AI optimizes for functionality, not for security, unless you explicitly ask it to. This chapter gives you the tools to identify the most common problems (OWASP Top~10), but it \textbf{does not replace a professional security review}.

\textbf{Before a commercial deployment} --- an application that handles real customer data, financial transactions, health data, or any sensitive information --- have the code reviewed by a cybersecurity professional. The cost of an audit is a fraction of the reputational and legal damage of a data breach.

The techniques you learn here are your \textbf{first line of defense}, not the last.
\end{dangerbox}

% ---------------------------------------------------------
\section{Backend Testing}

\subsection{The principle: Tests as Documentation}

AI-generated tests have a double value: they verify behavior \textbf{and} document the expected behavior. If tomorrow you change something and a test fails, the error message tells you exactly what broke.

\begin{praticabox}[title={Practice --- Generate backend tests}]
\begin{lstlisting}[language={}]
Generate complete tests for the Notes API backend with vitest.

Setup:
1. Install vitest and supertest as devDependencies
2. Create vitest.config.js with a PostgreSQL test database
3. Run migrations before the tests, clean the data between tests

Required tests:

test/unit/
  - note-service.test.js: note service CRUD
    (create, read, update, delete, note not found, access denied)
  - auth-service.test.js: token tests (JWT generation, verification,
    refresh, expired token)
  - validation.test.js: Zod middleware (empty title,
    content too long, extra fields ignored)

test/integration/
  - notes-api.test.js: endpoints with supertest
    (GET /api/notes without auth -> 401,
     POST with valid data -> 201,
     GET /:id of another user -> 403,
     DELETE /:id -> 200 + note removed)
  - auth-api.test.js: auth flow
    (GET /api/auth/me without token -> 401,
     with valid token -> 200 + profile,
     POST /api/auth/refresh with valid refresh -> new access)

Each test must have a descriptive name that explains the scenario.
\end{lstlisting}
\end{praticabox}

\begin{lstlisting}[language=bash,title={Run the tests}]
cd backend
npx vitest run
\end{lstlisting}

Expected output:
\begin{lstlisting}[language={}]
 v note-service.test.js (6 tests)
 v auth-service.test.js (4 tests)
 v validation.test.js (3 tests)
 v notes-api.test.js (4 tests)
 v auth-api.test.js (3 tests)

 Test Files  5 passed
 Tests       20 passed
\end{lstlisting}

\begin{tipbox}
If a test fails, copy the error and paste it into Copilot: ``This test fails with this error. What's the problem?'' The AI usually identifies the bug on the first attempt.
\end{tipbox}

% ---------------------------------------------------------
\section{Frontend Testing}

\begin{praticabox}[title={Practice --- Generate React tests}]
\begin{lstlisting}[language={}]
Generate tests for the React frontend with vitest and @testing-library/react.

1. Install: vitest, @testing-library/react,
   @testing-library/jest-dom, jsdom, msw (Mock Service Worker)
2. Configure MSW to simulate the backend

Required tests:
test/components/
  - NoteCard.test.jsx: renders title and content
  - NoteForm.test.jsx: field validation, submit callback
  - ProtectedRoute.test.jsx: redirects if not authenticated

test/hooks/
  - useNotes.test.jsx: loads notes, handles errors, updates

test/pages/
  - LoginPage.test.jsx: shows OAuth buttons
  - DashboardPage.test.jsx: loading, notes list, empty state
\end{lstlisting}
\end{praticabox}

\begin{lstlisting}[language=bash]
cd frontend
npx vitest run
\end{lstlisting}

% ---------------------------------------------------------
\section{Flutter App Testing}

\begin{praticabox}[title={Practice --- Generate Flutter tests}]
\begin{lstlisting}[language={}]
Generate tests for the Flutter app:

test/
  unit/
    - token_service_test.dart: save, read, delete
    - note_model_test.dart: fromJson, toJson, required fields

  widget/
    - login_screen_test.dart: OAuth buttons, loading state
    - note_card_test.dart: renders title, tap callback
    - dashboard_screen_test.dart: loading -> notes list, empty state

Use mockito to simulate the services.
Use ProviderScope.overrides to inject mock providers.
\end{lstlisting}
\end{praticabox}

\begin{lstlisting}[language=bash]
cd notes_mobile
flutter test
\end{lstlisting}

% ---------------------------------------------------------
\section{Confidence Tagging}\index{Confidence Tagging}

Not all AI-generated code requires the same level of review attention. \textbf{Confidence Tagging} classifies code by risk level.

\subsection{The risk matrix}

\begin{center}
\begin{tabularx}{\textwidth}{l l X}
\toprule
\textbf{Level} & \textbf{What it includes} & \textbf{Action} \\
\midrule
\textcolor{green!60!black}{\textbf{LOW}} & UI, styling, stateless components, data models & Quick review: does the code work? \\
\textcolor{orange}{\textbf{MEDIUM}} & Business logic, validation, routing, state management & Careful review: is the logic correct? \\
\textcolor{red}{\textbf{HIGH}} & Auth, authorization, crypto, DB queries, user input & Rigorous review: is it secure? \\
\bottomrule
\end{tabularx}
\end{center}

\begin{praticabox}[title={Practice --- Project risk analysis}]
\begin{lstlisting}[language={}]
Analyze the entire project and classify each file:

- LOW RISK: they do not handle sensitive data or critical logic
- MEDIUM RISK: business logic or validation
- HIGH RISK: authentication, authorization, user data, DB queries

For each HIGH RISK file, list the vulnerabilities to verify:

HIGH: backend/src/middleware/auth.js
   - JWT: verify that the secret is not hardcoded
   - Token: make sure it handles malformed tokens
   - User lookup: not vulnerable to injection

MEDIUM: backend/src/services/note-service.js
   - Verify that it always filters by userId

LOW: frontend/src/components/ui/Button.jsx
   - No checks needed
\end{lstlisting}
\end{praticabox}

% ---------------------------------------------------------
\section{Code Review of AI-Generated Code}\index{code review}

Confidence Tagging tells you \textbf{what} to review. This section teaches you \textbf{how} to do it --- even if you are not an experienced developer.

\subsection{Why AI gets things wrong differently from humans}

A human programmer makes distraction mistakes: typos, forgotten parentheses, variables with the wrong names. AI almost never makes these mistakes. AI bugs are more insidious:

\begin{center}
\begin{small}
\begin{tabularx}{\textwidth}{l l X}
\toprule
\textbf{Type of bug} & \textbf{Example} & \textbf{Why it happens} \\
\midrule
Happy path only & Works with valid input, crashes with empty input & AI optimizes for ``it works'', not for ``it is resilient'' \\
Non-existent APIs & Uses a method that does not exist in the installed version & It mixed documentation from different versions \\
Plausible but wrong logic & A filter that looks correct but excludes edge cases & The pattern looked right in the training data \\
Forgotten security & Endpoint without an authorization check & You did not ask for it explicitly \\
Phantom dependencies & Imports a package not in package.json & It saw it in a similar project \\
\bottomrule
\end{tabularx}
\end{small}
\end{center}

\subsection{Review methodology by level}

\textbf{\textcolor{green!60!black}{LOW RISK} --- Functional verification (2 minutes/file)}

UI components, styling, data models. The question is simple: \textit{does it work?}

\begin{enumerate}
  \item Did the AI create the file? Open it and verify that it is not empty
  \item Start the app and see whether the component is displayed
  \item If yes, move on
\end{enumerate}

\textbf{\textcolor{orange}{MEDIUM RISK} --- Logic verification (5--10 minutes/file)}

Business logic, validation, routing. The question: \textit{does it do what it should?}

\begin{enumerate}
  \item Read the main function and ask yourself: ``What happens if the input is empty? Null? Huge?''
  \item Look for dangerous keywords: \texttt{any}, \texttt{as any} (TypeScript), \texttt{!} (force unwrap), empty \texttt{catch} blocks
  \item Verify that the Zod/Joi validations match the requirements in \texttt{\_CONTEXT.md}
  \item If the logic seems correct but you are not sure, ask the AI: \textit{``Explain step by step what this function does and which edge cases it might not handle''}
\end{enumerate}

\textbf{\textcolor{red}{HIGH RISK} --- Rigorous verification (15--30 minutes/file)}

Authentication, authorization, crypto, DB queries. The question: \textit{is it secure?}

\begin{enumerate}
  \item \textbf{Authorization}: does every endpoint that accesses user data filter by \texttt{userId}? Can one user access another user's data?
  \item \textbf{JWT}: is the secret in an environment variable? Does the token have an expiration?
  \item \textbf{User input}: does all input go through validation before reaching the database?
  \item \textbf{Errors}: do production error messages \textsc{not} expose stack traces or internal details?
  \item Use the Security Checklist as a systematic guide
\end{enumerate}

\begin{praticabox}[title={Practice --- AI-assisted review}]
\begin{lstlisting}[language={}]
Analyze the HIGH RISK files of the notes-fullstack project.

For each file:
1. List all possible unhandled edge cases
2. Identify every point where user input reaches the database
   without going through validation
3. Verify that there are no hardcoded secrets
4. Check that every protected endpoint verifies authorization

Output format for each file:
VERIFIED: [what is good]
WARNING: [what could be a problem]
BUG: [what is definitely wrong, with fix]
\end{lstlisting}
\end{praticabox}

\begin{tipbox}
\textbf{The golden rule of review}: You do not need to understand \textit{how} every line of code works. You need to understand \textit{what} the function does and ask yourself: ``What could go wrong?''. The AI is your microscope --- you need to know where to point it.
\end{tipbox}

% ---------------------------------------------------------
\section{OWASP Top 10 Analysis}\index{OWASP!Top 10}

The 10 most common vulnerabilities in web applications~\cite{owasp:top10}. Let us verify them in our project.

\begin{praticabox}[title={Practice --- OWASP audit}]
\begin{lstlisting}[language={}]
Run a security analysis based on OWASP Top 10 (2021).

A01: Broken Access Control
- Can one user access another user's notes?
- Are admin endpoints protected?
- Do protected routes ALWAYS verify the JWT?

A02: Cryptographic Failures
- Is the JWT secret sufficiently long?
- Are passwords hashed with bcrypt?
- Does the database use SSL in production?

A03: Injection
- Are Prisma queries parameterized?
- Is the input validated with Zod before the database?

A04: Insecure Design
- Does the refresh token have an expiration?
- Is there rate limiting on login?

A05: Security Misconfiguration
- Do production errors expose internal details?
- Is Helmet.js configured?
- Is CORS configured for production?

A06: Vulnerable Components
- Does npm audit report critical vulnerabilities?

A07: Authentication Failures
- Do JWT tokens expire?
- Does logout invalidate the refresh token in the database?

A08: Software and Data Integrity
- Is package-lock.json committed?

A09: Security Logging
- Are failed login attempts logged?
- Do the logs NOT contain sensitive data?

A10: SSRF
- Does the backend make requests to URLs provided by the user?

For each vulnerability found, provide the fix with the code.
\end{lstlisting}
\end{praticabox}

\subsection{Common fixes}

\textbf{Rate Limiting:}\index{rate limiting}
\begin{lstlisting}[language={}]
Add rate limiting to the backend:
1. Install express-rate-limit
2. Limit /api/auth/* to 10 requests/minute per IP
3. Limit /api/notes to 100 requests/minute per user
\end{lstlisting}

\textbf{Helmet.js:}\index{Helmet.js}
\begin{lstlisting}[language={}]
Add Helmet.js for security headers:
1. Install helmet
2. Add helmet() as the first middleware
3. Configure CSP for production
\end{lstlisting}

\textbf{Dependency audit:}
\begin{lstlisting}[language=bash]
cd backend && npm audit
cd frontend && npm audit
npm audit fix  # if vulnerabilities are reported
\end{lstlisting}

% ---------------------------------------------------------
\section{Final Security Checklist}

\begin{praticabox}[title={Practice --- Full verification}]
\begin{center}
\begin{small}
\begin{tabularx}{\textwidth}{c l X}
\toprule
\textbf{\#} & \textbf{Check} & \textbf{Status} \\
\midrule
1 & JWT secret $\geq$ 256 bit, in environment variable & $\square$ \\
2 & Access token expires in 1 hour & $\square$ \\
3 & Refresh token expires in 7 days & $\square$ \\
4 & Logout deletes the refresh token from the database & $\square$ \\
5 & Notes filtered by userId in \textsc{all} queries & $\square$ \\
6 & Input validated with Zod before the database & $\square$ \\
7 & Production errors \textsc{do not} expose stack traces & $\square$ \\
8 & Helmet.js active with security headers & $\square$ \\
9 & Rate limiting on auth endpoints & $\square$ \\
10 & CORS configured only for allowed domains & $\square$ \\
11 & npm audit without critical vulnerabilities & $\square$ \\
12 & Mobile token in flutter\_secure\_storage & $\square$ \\
13 & Frontend does not save tokens in localStorage & $\square$ \\
14 & Database uses SSL in production & $\square$ \\
15 & \texttt{.env} and \texttt{keystore.jks} in \texttt{.gitignore} & $\square$ \\
\bottomrule
\end{tabularx}
\end{small}
\end{center}
\end{praticabox}

\begin{checkpointbox}
If all 15 checks are ticked and the tests pass, the application is ready for production deployment.
\end{checkpointbox}

% ---------------------------------------------------------
\section{Agentic Security: Beyond OWASP Web}

The OWASP Top~10 analysis protects the \textbf{application} you build. But who protects the \textbf{development infrastructure}? When you grant an AI agent access to databases, the file system, and APIs through MCP, you introduce new threat vectors that traditional OWASP does not cover.

In 2025, OWASP published the \textbf{OWASP Top~10 for MCP}\index{OWASP!MCP Top 10}~\cite{owasp:mcp-top10,owasp:ai-agent}, a classification specific to protocol vulnerabilities. Three of these deserve immediate attention:

\subsection{Tool Poisoning (MCP03)}\index{Tool Poisoning}\index{MCP!Tool Poisoning}

When you configure an MCP server, the agent trusts the \textbf{descriptions} of the tools exposed by the server to decide which action to execute. A compromised server can insert malicious descriptions that induce the agent to execute destructive operations.

\begin{dangerbox}
A tampered PostgreSQL MCP server could describe the tool \texttt{cleanup\_old\_data} as ``removes test records'', when in reality it executes \texttt{DROP TABLE users CASCADE}. The agent, trusting the description, would invoke it without hesitation.
\end{dangerbox}

\textbf{Mitigation:} install \textit{only} MCP servers from verified repositories (official \texttt{@modelcontextprotocol/} packages or projects with thousands of GitHub stars). Review the server source code before adopting it.

\subsection{Shadow MCP Servers (MCP09)}\index{Shadow MCP Server}

As happened with ``Shadow IT'', developers may install unverified MCP servers to speed up integration with new APIs. These servers bypass the team's security controls and create hidden entry points.

\textbf{Mitigation:} keep a shared registry of approved MCP servers (in the project \texttt{\_CONTEXT.md} or in team policies).

\subsection{Context Injection and Over-Sharing (MCP10)}\index{Context Injection}\index{MCP!Context Injection}

When the context window is shared across different sessions, sensitive information --- API tokens, secrets, user data --- can leak from one task to another.

\textbf{Mitigation:} always use \texttt{.env} files for secrets, \textbf{never} inline in MCP configurations. Verify that \texttt{.gitignore} includes all sensitive files. In team environments, consider mutual~TLS to authenticate MCP connections.

\begin{praticabox}[title={Practice --- Agentic security audit}]
\begin{lstlisting}[language={}]
Analyze the security of our agentic infrastructure:

1. TOOL VERIFICATION:
   - List all configured MCP servers (.vscode/mcp.json)
   - For each one: verify the source, check GitHub stars
     and the date of the last commit
   - Flag unofficial or unmaintained servers

2. SECRET MANAGEMENT:
   - Verify that NO secret is hardcoded in mcp.json
   - All secrets must be in environment variables
   - Is the .env file in .gitignore?

3. CONTEXT HYGIENE:
   - Does _CONTEXT.md contain credentials?
   - Do session logs contain sensitive data?
   - Do MCP configurations expose internal URLs or tokens?
\end{lstlisting}
\end{praticabox}

% ---------------------------------------------------------
\section{Testing Non-Deterministic Systems}

The unit and integration tests you wrote in this chapter validate \textbf{deterministic code}: given an input, the output is always the same. But in 2026 many applications integrate AI agents that produce \textbf{probabilistic} responses --- chatbots, recommendation engines, autonomous reasoning pipelines.

In these scenarios, binary pass/fail assertions break down: the ``correct'' output has many possible forms.

\subsection{Strategies for probabilistic testing}

\begin{center}
\begin{tabularx}{\textwidth}{l X}
\toprule
\textbf{Technique} & \textbf{When to use it} \\
\midrule
Semantic evaluation & The output must \textit{mean} the same thing, not be identical letter by letter \\
Golden dataset & Compare the responses with a reference set curated by human experts \\
Evaluator Agent & A second agent evaluates the quality of the first agent's output \\
Observability (tracing) & Trace the agent's \textit{decision tree} to understand \textit{why} it gave that response \\
Probabilistic fuzzing & Inject variations into the input to verify the robustness of the output \\
\bottomrule
\end{tabularx}
\end{center}

\begin{tipbox}
Platforms such as \textbf{Maxim~AI}~\cite{maxim:eval-2026} make it possible to evaluate multi-agent systems by measuring reasoning accuracy across multi-turn interactions. Open-source tools such as \textbf{Arize Phoenix}~\cite{arize:phoenix} (based on OpenTelemetry) trace the entire execution tree of an agent's decisions.
\end{tipbox}

\begin{notebox}
For the applications in this book (CRUD, REST~API, authentication), traditional tests are perfectly adequate. Probabilistic testing becomes necessary only when you integrate AI models \textit{inside} your application --- for example, a chatbot or a semantic search engine.
\end{notebox}

% ---------------------------------------------------------
\section*{Summary}

\begin{center}
\begin{tabularx}{\textwidth}{l X}
\toprule
\textbf{Aspect} & \textbf{Detail} \\
\midrule
Backend Tests & vitest + supertest, unit + integration \\
Frontend Tests & vitest + testing-library + MSW \\
Mobile Tests & flutter test + mockito \\
Risk & 3-level Confidence Tagging \\
Code Review & Review methodology by risk level \\
Web Security & OWASP Top~10 audit with fixes \\
Agentic Security & OWASP MCP Top~10: Tool Poisoning, Shadow MCP, Context Injection \\
Probabilistic Testing & Strategies for non-deterministic output \\
Checklist & 15 verifiable security checks \\
\bottomrule
\end{tabularx}
\end{center}
